\section*{Free Gaussian Wave Packets in 2D}

For a narrow scalar wave packet centered around momentum $\mathbf k_0$,
\[
\phi(t,\mathbf r)\approx
\mathcal A(t)
\exp\!\left[
-\frac{\Delta_\parallel^2}{2\sigma_\parallel^2(t)}
-\frac{\Delta_\perp^2}{2\sigma_\perp^2(t)}
\right]
\cos\Phi(t,\mathbf r),
\qquad
\mathbf r=(x,y).
\]

The packet center follows the group velocity
\[
\mathbf v_g=\frac{\mathbf k_0}{\omega_0},
\qquad
\omega_0=\sqrt{k_0^2+m^2}.
\]

The longitudinal and transverse widths broaden as
\[
\sigma_\parallel^2(t)=\sigma^2+\frac{\alpha_\parallel^2 t^2}{\sigma^2},
\qquad
\sigma_\perp^2(t)=\sigma^2+\frac{\alpha_\perp^2 t^2}{\sigma^2},
\]
with
\[
\alpha_\parallel=\frac{m^2}{\omega_0^3},
\qquad
\alpha_\perp=\frac{1}{\omega_0}.
\]

The amplitude decays due to spreading,
\[
\mathcal A(t)=
A\frac{\sigma^2}{
(\sigma^4+\alpha_\parallel^2 t^2)^{1/4}
(\sigma^4+\alpha_\perp^2 t^2)^{1/4}}.
\]

For two free packets,
\[
\phi(t,\mathbf r)=\phi_1(t,\mathbf r)+\phi_2(t,\mathbf r),
\]
so the overlap pattern is pure linear interference. No interaction term is present.
